\documentclass[11pt]{article}

\usepackage[utf8]{inputenc}
\usepackage[T1]{fontenc}
\usepackage{lmodern}
\usepackage{geometry}
\usepackage{amsmath,amssymb,amsfonts,amsthm,mathtools}
\usepackage{bm}
\usepackage{enumitem}
\usepackage{microtype}
\usepackage{hyperref}
\usepackage{titlesec}
\usepackage{booktabs}
\usepackage{array}
\usepackage{setspace}
\usepackage{mathrsfs}
\usepackage{physics}

\geometry{margin=1in}
\hypersetup{colorlinks=true, linkcolor=black, urlcolor=blue, citecolor=black}
\setlist[enumerate]{topsep=0.4em,itemsep=0.35em}
\setlist[itemize]{topsep=0.3em,itemsep=0.25em}
\titleformat{\section}{\large\bfseries}{\thesection.}{0.5em}{}
\titleformat{\subsection}{\normalsize\bfseries}{\thesubsection.}{0.5em}{}
\setstretch{1.06}

\newcommand{\Aether}{\AE{}ther}
\newcommand{\Aflow}{\AE{}ther-flow}
\newcommand{\TheoryName}{\Aflow{} Interpretation of Relativity}
\newcommand{\TheoryProgram}{\Aether{} / \Aflow{} framework}

\newcommand{\CompletedMetric}{g^{\sharp}}

\newcommand{\Car}{\mathrm{car}}
\newcommand{\Cur}{\mathrm{cur}}
\newcommand{\Hom}{\mathrm{hom}}

\newcommand{\ForSpace}[1]{\mathcal I_{#1}^{\mathrm{for}}}
\newcommand{\PrimShell}{\mathcal S_{\mathrm{prim}}}

\newcommand{\Reduction}[1]{\mathcal R_{#1}}
\newcommand{\CurrentCarrier}[1]{\mathfrak C_{#1}^{\mathrm{cur}}}
\newcommand{\EqCarrier}[1]{\mathfrak C_{#1}^{\mathrm{eq}}}
\newcommand{\CurrentExtract}[1]{\mathscr E_{#1}^{\mathrm{cur}}}
\newcommand{\EqExtract}[1]{\mathscr E_{#1}^{\mathrm{eq}}}
\newcommand{\PrimCurrent}[1]{\mathcal J_{#1}^{\mathrm{prim}}}
\newcommand{\PrimTheta}[1]{\Theta_{#1}^{\mathrm{prim}}}

\newcommand{\MetricOut}{\mathscr G_{\mathrm{out}}}
\newcommand{\MatterOut}{\mathscr T_{\mathrm{out}}}
\newcommand{\DeepMetricOut}{\mathscr G_{\mathrm{deep}}}
\newcommand{\DeepMatterOut}{\mathscr T_{\mathrm{deep}}}

\newcommand{\VisibleAffine}[1]{\widehat X_{#1}}
\newcommand{\DataSpace}[1]{\mathcal Y_{#1}^{\mathrm{obs}}}
\newcommand{\AmbientTarget}[1]{E_{#1}^{\mathrm{amb}}}

\newcommand{\ProtoSpace}[1]{\mathcal M_{#1}}
\newcommand{\ProtoBody}[1]{\mathcal P_{#1}}
\newcommand{\ResponseSpace}[1]{\mathcal R_{#1}}
\newcommand{\ResponseMap}[1]{\Xi_{#1}}
\newcommand{\ResponseData}[1]{\Xi_{#1}^{\mathrm{dat}}}
\newcommand{\ResponseShell}[1]{\Xi_{#1}^{\mathrm{sh}}}
\newcommand{\ResponseHidden}[1]{\Xi_{#1}^{\mathrm{hid}}}
\newcommand{\ResponseImage}[1]{\mathcal B_{#1}^{\mathrm{resp}}}
\newcommand{\GapSpace}[1]{H_{#1}}
\newcommand{\GapBody}[1]{D_{#1}}
\newcommand{\HiddenOperator}[1]{K_{#1}^{\mathrm{hid}}}
\newcommand{\ResponseSym}[1]{\mathbb H_{#1}}
\newcommand{\ResponseTime}[1]{\rho_{#1}}
\newcommand{\ProtoEmbed}[1]{\zeta_{#1}}

\newcommand{\CoreState}[1]{\mathcal C_{#1}}
\newcommand{\CoreSpace}[1]{\mathcal Q_{#1}}
\newcommand{\CoreAffine}[1]{\mathcal F_{#1}^{\mathrm{co}}}
\newcommand{\CoreBody}[1]{Q_{#1}^{\mathrm{co}}}
\newcommand{\CoreStateBody}[1]{P_{#1}^{\mathrm{co}}}
\newcommand{\CoreStateProj}[1]{\Pi_{#1}^{\mathrm{co,st}}}
\newcommand{\CoreDataProj}[1]{\Pi_{#1}^{\mathrm{co,dat}}}
\newcommand{\CoreShellProj}[1]{\Pi_{#1}^{\mathrm{co,sh}}}
\newcommand{\CoreDataMap}[1]{\Omega_{#1}^{\mathrm{co,dat}}}
\newcommand{\CoreShellMap}[1]{\Omega_{#1}^{\mathrm{co,sh}}}
\newcommand{\CoreSym}[1]{\mathbb K_{#1}^{\mathrm{co}}}
\newcommand{\CoreTime}[1]{\tau_{#1}^{\mathrm{co}}}
\newcommand{\CoreEmbed}[1]{\eta_{#1}}

\newcommand{\GraphSpace}[1]{\mathcal G_{#1}}
\newcommand{\GraphAffine}[1]{\mathcal F_{#1}^{\mathrm{gr}}}
\newcommand{\GraphBody}[1]{Q_{#1}^{\mathrm{gr}}}
\newcommand{\StateProj}[1]{\Pi_{#1}^{\mathrm{st}}}
\newcommand{\GraphDataProj}[1]{\Pi_{#1}^{\mathrm{gr,dat}}}
\newcommand{\GraphShellProj}[1]{\Pi_{#1}^{\mathrm{gr,sh}}}
\newcommand{\GraphEmbed}[1]{\upsilon_{#1}}

\newcommand{\DeepSpace}[1]{\mathcal U_{#1}}
\newcommand{\ChainSelect}[1]{\mathcal S_{#1}}
\newcommand{\DeepShell}{\mathcal U_{\mathrm{tot}}}
\newcommand{\ChainSelectTriple}{\mathcal S_{\mathrm{tot}}}

\newtheorem{definition}{Definition}
\newtheorem{lemma}{Lemma}
\newtheorem{proposition}{Proposition}
\newtheorem{remark}{Remark}
\newtheorem{corollary}{Corollary}
\newtheorem{theorem}{Theorem}

\title{\textbf{The \TheoryName{}}\\[0.4em]
\large Primitive-Reservoir Observer-Reduction\\
\large Affine Response-Graph Dynamics and Combined Response Substrate Law\\
\large Collapse to a Minimal Zero-Response Affine Completion Core}
\author{}
\date{}

\begin{document}

\maketitle

\begin{abstract}
The current deepest benchmark-facing note on the active primitive-reservoir observer-reduction line already sharpened the remaining burden once more: the combined response substrate law is no longer merely available, because it is derived canonically from affine response-graph dynamics and is necessary there up to the forced graph-state and factor-preserving response-coordinate identifications. The remaining honest burden is therefore no longer to justify the combined response law. It is to explain why the affine response-graph layer itself should not become the new stopping point.

The present note makes the strongest honest move now available without overstating first-principles closure. Rather than insert a looser deeper same-output relay, it proves that the affine response-graph layer and the combined-response layer both collapse to a smaller zero-response affine completion core with hidden-sector gap. For each sector this core consists of one affine completion plane in completion-state, observer-data, and shell-response coordinates; one compact convex completion body in that plane; one compact convex hidden body with a positive hidden-sector gap operator; symmetry and time-reversal actions preserving the core data; an exact zero-response shell realization; fixed-data solvability on the zero-response slice; and a core-kernel coercivity inequality on data-invisible variations.

From that core one derives canonically the compact-convex primitive substrate body, the single affine combined response map to observer data, shell response, and hidden response, the response-image factorization, the exact shell realization, the fixed-data solvability property, and the response-kernel coercivity mechanism. One also derives the affine graph plane, the compact convex graph body, the graph-state projection producing the primitive body, the unique affine graph function, the graph-level hidden splitting with positive gap operator, the symmetry / time-reversal structure, the exact zero-response shell realization, the fixed-data solvability property, and the graph-kernel coercivity mechanism. Thus the graph package carries no extra operative freedom beyond the completion core.

The benchmark package remains fixed throughout. The unique selector determined by the zero-response shell locus still pulls back exactly the same one operative metric \(\CompletedMetric\) and exactly the same ordinary-matter route through that metric, with no second operative metric and no enlarged low-energy matter law. This remains a disciplined derivational gain rather than a completed final microphysical derivation of GR.
\end{abstract}

\section{Introduction}

The active derivational program of \TheoryName{} has again narrowed what counts as an honest next step. The observer-reduction datum line, the defect-fiber factorization theorem, the one-metric output theorem, the chain-forcing theorem, the selector-existence theorem, the stationary-construction theorem, the explicit-package theorem, the primitive-normal-form theorem, the ambient-shell-law theorem, and the later combined-response theorem have each discharged what was honestly open at their respective layers. The most recent of those notes matters because it changes what now counts as genuine upstream progress. Once the combined response substrate law is itself forced by affine response-graph dynamics, another note that merely restates that law no longer counts as the front-line move.

The remaining honest burden named by the previous note is narrower. It is not to justify the combined response law again, but to justify why the affine response-graph dynamics package itself is forced on the active line rather than becoming the new disciplined stopping point. There are, in principle, two ways to make that move honestly: derive the graph package from still more primitive substrate dynamics, or prove a sharper inevitability theorem that removes the remaining graph-level freedom. The present note takes the second route.

Its gain is stronger than another deeper same-output relay and narrower than a completed ultimate microphysical derivation. The key observation is that the graph layer still records redundant presentation data. What actually matters for the downstream benchmark package is a smaller zero-response affine completion core together with a positively gapped hidden factor. Once that core is fixed, the combined response law and the graph-style presentation are both forced. The graph layer therefore ceases to be a genuine new freedom.

This note does not alter the benchmark package. It does not change the one operative completed metric, it does not change the ordinary matter route, it does not introduce a second operative metric, and it does not enlarge the low-energy matter law. Its gain is structural and upstream: the affine response-graph layer is shown not to be a new deepest package, but a redundant presentation of a smaller completion core.

\paragraph{Framework and claim boundary.}
\TheoryProgram{} denotes the broader \Aether{} / \Aflow{} framework built on the claims that the \Aether{} is the underlying four-dimensional substrate of reality, the \Aflow{} is its intrinsic ordered motion, observed three-dimensional space is the local experiential slice of that deeper substrate, S-time is the experienced order of change arising from matter, light, and the \Aflow{}, and observed expansion is the three-dimensional appearance of deeper four-dimensional ordered motion. Within that framework, \TheoryName{} denotes the exact-closure theory in which the effective gravitational dynamics are adopted to be exactly Einsteinian with universal matter coupling. In the active sequence, \emph{adoption} means use of that established relativistic dynamics without claiming substrate derivation, while \emph{derivation} is reserved for a first-principles recovery from explicit substrate variables.

The present note stays strictly inside that boundary. It does not claim a completed ultimate microphysical theory. Its narrower theorem is only that, on the active primitive-observer-reduction line, the affine response-graph layer and the combined-response layer both collapse to a minimal zero-response affine completion core with hidden-sector gap.

\section{Fixed Primitive Continuation Data}

\begin{definition}[Recorded primitive continuation data]
\label{def:fixed}
Fix the following continuation data from the active primitive-reservoir observer-reduction line.
\begin{enumerate}
    \item For each sector label \(\sigma\in\{\Car,\Cur,\Hom\}\), a primitive forcing shell
    \[
    \ForSpace{\sigma},
    \]
    a visible reduction map
    \[
    \Reduction{\sigma}:\ForSpace{\sigma}\to X_{\sigma},
    \]
    and shell-level current and equilibrium carriers
    \[
    \CurrentCarrier{\sigma}:\ForSpace{\sigma}\to Z_{\sigma}^{\mathrm{cur}},
    \qquad
    \EqCarrier{\sigma}:\ForSpace{\sigma}\to Z_{\sigma}^{\mathrm{eq}}.
    \]
    \item The product shell
    \[
    \PrimShell
    \equiv
    \ForSpace{\Car}\times \ForSpace{\Cur}\times \ForSpace{\Hom}.
    \]
    \item Fixed shell extractors
    \[
    \CurrentExtract{\sigma}:Z_{\sigma}^{\mathrm{cur}}\to Y_{\sigma}^{\mathrm{cur}},
    \qquad
    \EqExtract{\sigma}:Z_{\sigma}^{\mathrm{eq}}\to Y_{\sigma}^{\mathrm{eq}},
    \]
    together with the recorded shell composites
    \begin{equation}
    \PrimCurrent{\sigma}
    =
    \CurrentExtract{\sigma}\circ \CurrentCarrier{\sigma},
    \qquad
    \PrimTheta{\sigma}
    =
    \EqExtract{\sigma}\circ \EqCarrier{\sigma}.
    \label{eq:primcomposites}
    \end{equation}
    \item The recorded metric and ordinary-matter outputs
    \[
    \MetricOut:\PrimShell\to\{\text{observer metrics}\},
    \qquad
    \MatterOut:\PrimShell\times\Psi\to\{\text{observer stress tensors}\},
    \]
    satisfying
    \begin{equation}
    \MetricOut(\mathbf w)=\CompletedMetric,
    \qquad
    \MatterOut(\mathbf w,\psi)
    =
    T_{\mu\nu}^{\mathrm{matter}}[\CompletedMetric,\psi]
    \label{eq:fixedoutputs}
    \end{equation}
    for every \(\mathbf w\in\PrimShell\) and every matter configuration \(\psi\in\Psi\).
\end{enumerate}
\end{definition}

\begin{remark}
Definition \ref{def:fixed} is not reopened here. The primitive shell data and the benchmark one-metric / ordinary-matter outputs remain fixed continuation data. The new burden begins one level deeper again: why the affine response-graph package of the previous note should not itself remain the deepest disciplined assumption.
\end{remark}

\section{Target Combined Response Law}

\begin{definition}[Combined response substrate law with hidden-sector gap]
\label{def:combined}
Fix \(\sigma\in\{\Car,\Cur,\Hom\}\). A \emph{combined response substrate law with hidden-sector gap} for sector \(\sigma\) consists of the following data.
\begin{enumerate}
    \item A finite-dimensional Euclidean primitive substrate space
    \[
    \ProtoSpace{\sigma}
    \]
    with a compact convex body
    \[
    \ProtoBody{\sigma}\subset \ProtoSpace{\sigma}.
    \]
    \item An observer-data target space
    \[
    \DataSpace{\sigma}
    \equiv
    \VisibleAffine{\sigma}\times Z_{\sigma}^{\mathrm{cur}}\times Z_{\sigma}^{\mathrm{eq}},
    \]
    an ambient shell target
    \[
    \AmbientTarget{\sigma},
    \]
    a hidden response space
    \[
    \GapSpace{\sigma},
    \]
    and the product response space
    \[
    \ResponseSpace{\sigma}
    \equiv
    \DataSpace{\sigma}\oplus \AmbientTarget{\sigma}\oplus \GapSpace{\sigma}.
    \]
    \item An affine combined response map
    \[
    \ResponseMap{\sigma}
    =
    \bigl(
    \ResponseData{\sigma},
    \ResponseShell{\sigma},
    \ResponseHidden{\sigma}
    \bigr):
    \ProtoSpace{\sigma}\to \ResponseSpace{\sigma}.
    \]
    \item The response image body
    \begin{equation}
    \ResponseImage{\sigma}
    \equiv
    \ResponseMap{\sigma}(\ProtoBody{\sigma})
    \label{eq:imagebody}
    \end{equation}
    is compact and convex, and its shell slice
    \[
    \ResponseImage{\sigma}\cap
    \bigl(
    \DataSpace{\sigma}\oplus \AmbientTarget{\sigma}\oplus\{0\}
    \bigr)
    \]
    has nonempty relative interior in
    \[
    \DataSpace{\sigma}\oplus \AmbientTarget{\sigma}\oplus\{0\}.
    \]
    \item A compact convex hidden body
    \[
    \GapBody{\sigma}\subset \GapSpace{\sigma}
    \]
    with
    \[
    0\in\operatorname{int}\GapBody{\sigma},
    \]
    such that the response image factorizes exactly as
    \begin{equation}
    \ResponseImage{\sigma}
    =
    \left\{
    (u,e,n):
    (u,e,0)\in\ResponseImage{\sigma},
    \quad
    n\in\GapBody{\sigma}
    \right\}.
    \label{eq:responsefactor}
    \end{equation}
    \item A positive self-adjoint hidden-sector operator
    \[
    \HiddenOperator{\sigma}:\GapSpace{\sigma}\to\GapSpace{\sigma}
    \]
    with lower bound
    \begin{equation}
    \ev{\eta,\HiddenOperator{\sigma}\eta}
    \ge
    \kappa_{\sigma}^{\mathrm{hid}}\,
    \|\eta\|^{2}
    \qquad
    \text{for all }\eta\in\GapSpace{\sigma}
    \label{eq:combinedgap}
    \end{equation}
    for some \(\kappa_{\sigma}^{\mathrm{hid}}>0\).
    \item A compact symmetry group
    \[
    \ResponseSym{\sigma}\subset
    O(\DataSpace{\sigma})\times O(\AmbientTarget{\sigma})\times O(\GapSpace{\sigma})
    \]
    and an orthogonal involution
    \[
    \ResponseTime{\sigma}\in
    O(\DataSpace{\sigma})\times O(\AmbientTarget{\sigma})\times O(\GapSpace{\sigma})
    \]
    preserving \(\ResponseImage{\sigma}\), the hidden-zero slice
    \[
    \DataSpace{\sigma}\oplus \AmbientTarget{\sigma}\oplus\{0\},
    \]
    the zero shell value \(0\in \AmbientTarget{\sigma}\), and \(\HiddenOperator{\sigma}\).
    \item A smooth embedding
    \[
    \jmath_{\sigma}:X_{\sigma}\hookrightarrow \VisibleAffine{\sigma}
    \]
    and a smooth shell realization map
    \[
    \ProtoEmbed{\sigma}:\ForSpace{\sigma}\hookrightarrow \ProtoSpace{\sigma}
    \]
    such that
    \[
    \ProtoEmbed{\sigma}(\ForSpace{\sigma})
    \subset
    \ProtoBody{\sigma}\cap
    \bigl(\ResponseShell{\sigma}\bigr)^{-1}(0)\cap
    \bigl(\ResponseHidden{\sigma}\bigr)^{-1}(0),
    \]
    the map \(\ResponseMap{\sigma}\circ\ProtoEmbed{\sigma}\) is a diffeomorphism onto
    \[
    \ResponseImage{\sigma}\cap
    \bigl(
    \DataSpace{\sigma}\oplus\{0\}\oplus\{0\}
    \bigr),
    \]
    and
    \begin{equation}
    \ResponseData{\sigma}\bigl(\ProtoEmbed{\sigma}(w)\bigr)
    =
    \bigl(
    \jmath_{\sigma}(\Reduction{\sigma}(w)),
    \CurrentCarrier{\sigma}(w),
    \EqCarrier{\sigma}(w)
    \bigr)
    \label{eq:combinedcompat}
    \end{equation}
    for every \(w\in\ForSpace{\sigma}\).
    \item Zero-response shell solvability on fixed data fibers:
    \begin{equation}
    \operatorname{pr}_{\DataSpace{\sigma}}\bigl(\ResponseImage{\sigma}\bigr)
    =
    \operatorname{pr}_{\DataSpace{\sigma}}
    \Bigl(
    \ResponseImage{\sigma}\cap
    \bigl(
    \DataSpace{\sigma}\oplus\{0\}\oplus\{0\}
    \bigr)
    \Bigr).
    \label{eq:combinedsolvability}
    \end{equation}
    \item Response-kernel coercivity on the primitive shell: there exists \(\kappa_{\sigma}^{\mathrm{resp}}>0\) such that for every
    \[
    w\in\ForSpace{\sigma},
    \qquad
    \delta m\in T_{\ProtoEmbed{\sigma}(w)}\ProtoSpace{\sigma}
    \]
    satisfying
    \[
    D\ResponseData{\sigma}\bigl(\ProtoEmbed{\sigma}(w)\bigr)\,\delta m=0,
    \]
    one has
    \begin{equation}
    \bigl\|
    D\ResponseShell{\sigma}\bigl(\ProtoEmbed{\sigma}(w)\bigr)\,\delta m
    \bigr\|^{2}
    +
    \ev{
    D\ResponseHidden{\sigma}\bigl(\ProtoEmbed{\sigma}(w)\bigr)\,\delta m,
    \HiddenOperator{\sigma}
    D\ResponseHidden{\sigma}\bigl(\ProtoEmbed{\sigma}(w)\bigr)\,\delta m
    }
    \ge
    \kappa_{\sigma}^{\mathrm{resp}}\,
    \|\delta m\|^{2}.
    \label{eq:combinedcoercive}
    \end{equation}
\end{enumerate}
\end{definition}

\begin{remark}
Definition \ref{def:combined} is the benchmark-facing target package fixed by the previous note. The purpose here is not to reinsert it as the next deeper assumption, but to show that it is already a canonical output of smaller completion-core data.
\end{remark}

\section{Minimal Zero-Response Affine Completion Core}

\begin{definition}[Minimal zero-response affine completion core with hidden-sector gap]
\label{def:core}
Fix \(\sigma\in\{\Car,\Cur,\Hom\}\). A \emph{minimal zero-response affine completion core with hidden-sector gap} for sector \(\sigma\) consists of the following data.
\begin{enumerate}
    \item A finite-dimensional Euclidean completion-state space
    \[
    \CoreState{\sigma},
    \]
    an observer-data space
    \[
    \DataSpace{\sigma}
    \equiv
    \VisibleAffine{\sigma}\times Z_{\sigma}^{\mathrm{cur}}\times Z_{\sigma}^{\mathrm{eq}},
    \]
    an ambient shell-response space
    \[
    \AmbientTarget{\sigma},
    \]
    a hidden space
    \[
    \GapSpace{\sigma},
    \]
    and the zero-hidden completion space
    \[
    \CoreSpace{\sigma}
    \equiv
    \CoreState{\sigma}\oplus\DataSpace{\sigma}\oplus \AmbientTarget{\sigma}.
    \]
    Let
    \[
    \CoreStateProj{\sigma}:\CoreSpace{\sigma}\to \CoreState{\sigma},
    \qquad
    \CoreDataProj{\sigma}:\CoreSpace{\sigma}\to \DataSpace{\sigma},
    \qquad
    \CoreShellProj{\sigma}:\CoreSpace{\sigma}\to \AmbientTarget{\sigma}
    \]
    denote the coordinate projections.
    \item An affine completion plane
    \[
    \CoreAffine{\sigma}\subset \CoreSpace{\sigma}
    \]
    such that the restriction
    \[
    \CoreStateProj{\sigma}\big|_{\CoreAffine{\sigma}}:
    \CoreAffine{\sigma}\to \CoreState{\sigma}
    \]
    is an affine isomorphism.
    \item A compact convex completion body
    \[
    \CoreBody{\sigma}\subset \CoreAffine{\sigma}
    \]
    with nonempty relative interior in \(\CoreAffine{\sigma}\).
    \item A compact convex hidden body
    \[
    \GapBody{\sigma}\subset \GapSpace{\sigma}
    \]
    with
    \[
    0\in\operatorname{int}\GapBody{\sigma},
    \]
    together with a positive self-adjoint hidden-sector operator
    \[
    \HiddenOperator{\sigma}:\GapSpace{\sigma}\to\GapSpace{\sigma}
    \]
    satisfying
    \begin{equation}
    \ev{\eta,\HiddenOperator{\sigma}\eta}
    \ge
    \kappa_{\sigma}^{\mathrm{hid}}\,
    \|\eta\|^{2}
    \qquad
    \text{for all }\eta\in\GapSpace{\sigma}
    \label{eq:coregap}
    \end{equation}
    for some \(\kappa_{\sigma}^{\mathrm{hid}}>0\).
    \item A compact symmetry group
    \[
    \CoreSym{\sigma}\subset
    O(\CoreState{\sigma})\times
    O(\DataSpace{\sigma})\times
    O(\AmbientTarget{\sigma})\times
    O(\GapSpace{\sigma})
    \]
    and an orthogonal involution
    \[
    \CoreTime{\sigma}\in
    O(\CoreState{\sigma})\times
    O(\DataSpace{\sigma})\times
    O(\AmbientTarget{\sigma})\times
    O(\GapSpace{\sigma})
    \]
    preserving \(\CoreAffine{\sigma}\), \(\CoreBody{\sigma}\), the zero-response slice
    \[
    \CoreState{\sigma}\oplus \DataSpace{\sigma}\oplus\{0\},
    \]
    the hidden body \(\GapBody{\sigma}\), and \(\HiddenOperator{\sigma}\).
    \item A smooth embedding
    \[
    \jmath_{\sigma}:X_{\sigma}\hookrightarrow \VisibleAffine{\sigma}
    \]
    and a smooth zero-response shell realization
    \[
    \CoreEmbed{\sigma}:\ForSpace{\sigma}\hookrightarrow \CoreSpace{\sigma}
    \]
    whose image is exactly
    \[
    \CoreBody{\sigma}\cap
    \bigl(\CoreShellProj{\sigma}\bigr)^{-1}(0),
    \]
    and such that
    \begin{equation}
    \CoreDataProj{\sigma}\bigl(\CoreEmbed{\sigma}(w)\bigr)
    =
    \bigl(
    \jmath_{\sigma}(\Reduction{\sigma}(w)),
    \CurrentCarrier{\sigma}(w),
    \EqCarrier{\sigma}(w)
    \bigr)
    \label{eq:corecompat}
    \end{equation}
    for every \(w\in\ForSpace{\sigma}\).
    \item Fixed-data solvability on the zero-response slice:
    \begin{equation}
    \CoreDataProj{\sigma}\bigl(\CoreBody{\sigma}\bigr)
    =
    \CoreDataProj{\sigma}
    \Bigl(
    \CoreBody{\sigma}\cap
    \bigl(\CoreShellProj{\sigma}\bigr)^{-1}(0)
    \Bigr).
    \label{eq:coresolvability}
    \end{equation}
    \item Core-kernel coercivity on the zero-response shell: there exists \(\kappa_{\sigma}^{\mathrm{co}}>0\) such that for every
    \[
    w\in\ForSpace{\sigma},
    \qquad
    \delta c\in T_{\CoreEmbed{\sigma}(w)}\CoreAffine{\sigma}
    \]
    satisfying
    \[
    \CoreDataProj{\sigma}(\delta c)=0,
    \]
    one has
    \begin{equation}
    \bigl\|
    \CoreShellProj{\sigma}(\delta c)
    \bigr\|^{2}
    \ge
    \kappa_{\sigma}^{\mathrm{co}}\,
    \bigl\|
    \CoreStateProj{\sigma}(\delta c)
    \bigr\|^{2}.
    \label{eq:corecoercive}
    \end{equation}
\end{enumerate}
\end{definition}

\begin{remark}
Definition \ref{def:core} is smaller than the graph package of the previous note in exactly the sense needed here. It retains only the zero-hidden affine completion plane and body, the separate hidden factor with positive gap, the symmetry / time-reversal structure, the exact shell realization, the solvability statement, and the coercive control on data-invisible variations. It does not separately stipulate a primitive substrate body, a full combined response map, or a graph-level thickening of the same information.
\end{remark}

\section{Canonical Combined Response Law from the Completion Core}

\begin{definition}[Canonical combined response law induced by the completion core]
\label{def:canonical}
Assume Definition \ref{def:core}. Define the canonical combined response substrate law in sector \(\sigma\) as follows.
\begin{enumerate}
    \item The completion-state body is
    \begin{equation}
    \CoreStateBody{\sigma}
    \equiv
    \CoreStateProj{\sigma}\bigl(\CoreBody{\sigma}\bigr).
    \label{eq:corestatebody}
    \end{equation}
    \item The unique affine zero-hidden completion maps are
    \begin{equation}
    \CoreDataMap{\sigma}
    \equiv
    \CoreDataProj{\sigma}\circ
    \bigl(
    \CoreStateProj{\sigma}\big|_{\CoreAffine{\sigma}}
    \bigr)^{-1},
    \qquad
    \CoreShellMap{\sigma}
    \equiv
    \CoreShellProj{\sigma}\circ
    \bigl(
    \CoreStateProj{\sigma}\big|_{\CoreAffine{\sigma}}
    \bigr)^{-1}.
    \label{eq:coremaps}
    \end{equation}
    \item The primitive substrate space is
    \[
    \ProtoSpace{\sigma}
    \equiv
    \CoreState{\sigma}\oplus \GapSpace{\sigma},
    \]
    and the primitive substrate body is the product body
    \begin{equation}
    \ProtoBody{\sigma}
    \equiv
    \CoreStateBody{\sigma}\times \GapBody{\sigma}.
    \label{eq:protobody}
    \end{equation}
    \item The response space is
    \[
    \ResponseSpace{\sigma}
    \equiv
    \DataSpace{\sigma}\oplus \AmbientTarget{\sigma}\oplus \GapSpace{\sigma}.
    \]
    \item The affine combined response map is
    \begin{equation}
    \ResponseMap{\sigma}(m,n)
    \equiv
    \bigl(
    \CoreDataMap{\sigma}(m),
    \CoreShellMap{\sigma}(m),
    n
    \bigr).
    \label{eq:responsemap}
    \end{equation}
    \item The response image body is
    \begin{equation}
    \ResponseImage{\sigma}
    \equiv
    \ResponseMap{\sigma}(\ProtoBody{\sigma}).
    \label{eq:responseimage}
    \end{equation}
    \item The symmetry group and time reversal on the response factors are
    \[
    \ResponseSym{\sigma}
    \equiv
    \left\{
    (g_{\mathrm{dat}},g_{\mathrm{sh}},g_{\mathrm{hid}}):
    (g_{\mathrm{st}},g_{\mathrm{dat}},g_{\mathrm{sh}},g_{\mathrm{hid}})
    \in \CoreSym{\sigma}
    \right\},
    \]
    \[
    \ResponseTime{\sigma}
    \equiv
    \bigl(
    \CoreTime{\sigma}\big|_{\DataSpace{\sigma}},
    \CoreTime{\sigma}\big|_{\AmbientTarget{\sigma}},
    \CoreTime{\sigma}\big|_{\GapSpace{\sigma}}
    \bigr).
    \]
    \item The shell realization map is
    \begin{equation}
    \ProtoEmbed{\sigma}(w)
    \equiv
    \bigl(
    \CoreStateProj{\sigma}(\CoreEmbed{\sigma}(w)),
    0
    \bigr).
    \label{eq:protoembed}
    \end{equation}
\end{enumerate}
\end{definition}

\begin{proposition}[The completion core induces a combined response substrate law]
\label{prop:combined}
Assume Definitions \ref{def:core} and \ref{def:canonical}. Then the canonical data of Definition \ref{def:canonical} satisfy every requirement of Definition \ref{def:combined}.
\end{proposition}

\begin{proof}
Because \(\CoreStateProj{\sigma}\big|_{\CoreAffine{\sigma}}\) is an affine isomorphism and \(\CoreBody{\sigma}\subset\CoreAffine{\sigma}\) is compact and convex, \(\CoreStateBody{\sigma}\) is compact and convex with nonempty interior in \(\CoreState{\sigma}\). Hence \(\ProtoBody{\sigma}=\CoreStateBody{\sigma}\times\GapBody{\sigma}\) is compact and convex in \(\ProtoSpace{\sigma}\). Equation \eqref{eq:responsemap} is affine, so \eqref{eq:imagebody} holds by \eqref{eq:responseimage}.

Because \(\ResponseMap{\sigma}(m,n)=\bigl(\CoreDataMap{\sigma}(m),\CoreShellMap{\sigma}(m),n\bigr)\), the response image factorizes exactly as
\[
\ResponseImage{\sigma}
=
\left\{
(u,e,n):
(u,e,0)\in\ResponseImage{\sigma},
\quad
n\in\GapBody{\sigma}
\right\},
\]
which is \eqref{eq:responsefactor}. Equation \eqref{eq:combinedgap} is exactly \eqref{eq:coregap}. The symmetry and time-reversal requirements follow because every element of \(\CoreSym{\sigma}\) and \(\CoreTime{\sigma}\) preserves the completion body, the zero-response slice, the hidden body, and \(\HiddenOperator{\sigma}\).

For the shell realization, \eqref{eq:protoembed} maps \(\ForSpace{\sigma}\) into \(\ProtoBody{\sigma}\cap\bigl(\ResponseShell{\sigma}\bigr)^{-1}(0)\cap\bigl(\ResponseHidden{\sigma}\bigr)^{-1}(0)\). Since \(\CoreEmbed{\sigma}\) is a diffeomorphism onto \(\CoreBody{\sigma}\cap\bigl(\CoreShellProj{\sigma}\bigr)^{-1}(0)\) and \(\CoreStateProj{\sigma}\big|_{\CoreAffine{\sigma}}\) is an affine isomorphism, the map \(\ResponseMap{\sigma}\circ\ProtoEmbed{\sigma}\) is a diffeomorphism onto
\[
\ResponseImage{\sigma}\cap
\bigl(
\DataSpace{\sigma}\oplus\{0\}\oplus\{0\}
\bigr).
\]
Using \eqref{eq:corecompat},
\[
\ResponseData{\sigma}\bigl(\ProtoEmbed{\sigma}(w)\bigr)
=
\CoreDataProj{\sigma}\bigl(\CoreEmbed{\sigma}(w)\bigr)
=
\bigl(
\jmath_{\sigma}(\Reduction{\sigma}(w)),
\CurrentCarrier{\sigma}(w),
\EqCarrier{\sigma}(w)
\bigr),
\]
which is \eqref{eq:combinedcompat}.

The solvability condition \eqref{eq:combinedsolvability} follows directly from \eqref{eq:coresolvability}, because the hidden coordinate can be set to \(0\) independently. Finally, let
\[
w\in\ForSpace{\sigma},
\qquad
\delta m=(\delta m_{\mathrm{co}},\delta n)\in
T_{\ProtoEmbed{\sigma}(w)}\ProtoSpace{\sigma}
\]
with
\[
D\ResponseData{\sigma}\bigl(\ProtoEmbed{\sigma}(w)\bigr)\,\delta m=0.
\]
Then \(D\CoreDataMap{\sigma}\,\delta m_{\mathrm{co}}=0\). Let
\[
\delta c\in T_{\CoreEmbed{\sigma}(w)}\CoreAffine{\sigma}
\]
be the unique tangent vector satisfying
\[
\CoreStateProj{\sigma}(\delta c)=\delta m_{\mathrm{co}}.
\]
By \eqref{eq:coremaps},
\[
\CoreDataProj{\sigma}(\delta c)
=
D\CoreDataMap{\sigma}\,\delta m_{\mathrm{co}}
=
0.
\]
Applying \eqref{eq:corecoercive} gives
\[
\|\CoreShellProj{\sigma}(\delta c)\|^{2}
\ge
\kappa_{\sigma}^{\mathrm{co}}\,
\|\delta m_{\mathrm{co}}\|^{2}.
\]
Using \eqref{eq:coregap},
\[
\ev{\delta n,\HiddenOperator{\sigma}\delta n}
\ge
\kappa_{\sigma}^{\mathrm{hid}}\,
\|\delta n\|^{2}.
\]
Because
\[
\CoreShellProj{\sigma}(\delta c)
=
D\ResponseShell{\sigma}\bigl(\ProtoEmbed{\sigma}(w)\bigr)\,\delta m,
\qquad
\delta n
=
D\ResponseHidden{\sigma}\bigl(\ProtoEmbed{\sigma}(w)\bigr)\,\delta m,
\]
one obtains
\[
\bigl\|
D\ResponseShell{\sigma}\bigl(\ProtoEmbed{\sigma}(w)\bigr)\,\delta m
\bigr\|^{2}
+
\ev{
D\ResponseHidden{\sigma}\bigl(\ProtoEmbed{\sigma}(w)\bigr)\,\delta m,
\HiddenOperator{\sigma}
D\ResponseHidden{\sigma}\bigl(\ProtoEmbed{\sigma}(w)\bigr)\,\delta m
}
\ge
\kappa_{\sigma}^{\mathrm{resp}}\,
\|\delta m\|^{2}
\]
with \(\kappa_{\sigma}^{\mathrm{resp}}=\min\{\kappa_{\sigma}^{\mathrm{co}},\kappa_{\sigma}^{\mathrm{hid}}\}\). Thus \eqref{eq:combinedcoercive} holds.
\end{proof}

\begin{proposition}[The combined response package is forced by the completion core]
\label{prop:forcedcombined}
Assume Definition \ref{def:core}. Then:
\begin{enumerate}
    \item the compact-convex primitive substrate body is necessarily the product body \eqref{eq:protobody};
    \item the affine combined response map is necessarily the unique map \eqref{eq:responsemap} determined by the affine completion plane and the hidden factor;
    \item the response image is necessarily the factorized image \eqref{eq:responseimage};
    \item the hidden response factor and its positive operator are necessarily \(\GapBody{\sigma}\) and \(\HiddenOperator{\sigma}\);
    \item the symmetry / time-reversal structure, the exact zero-response shell realization, the fixed-data solvability statement, and the response-kernel coercivity mechanism are all fixed by the same completion core.
\end{enumerate}
\end{proposition}

\begin{proof}
Because \(\CoreStateProj{\sigma}\big|_{\CoreAffine{\sigma}}\) is an affine isomorphism, the zero-hidden completion plane is the graph of a unique affine pair \((\CoreDataMap{\sigma},\CoreShellMap{\sigma})\), namely \eqref{eq:coremaps}. Once the hidden body and operator are fixed, the only combined response law compatible with that completion plane is \eqref{eq:responsemap} on the product state space \(\CoreState{\sigma}\oplus\GapSpace{\sigma}\), and its body is necessarily \eqref{eq:protobody}. The remaining items follow from Definition \ref{def:canonical} and Proposition \ref{prop:combined}.
\end{proof}

\section{Compressed Graph Presentation and Collapse of Graph-Level Freedom}

\begin{definition}[Compressed affine response-graph presentation induced by the core]
\label{def:graph}
Assume Definitions \ref{def:core} and \ref{def:canonical}. Define the associated compressed affine response-graph presentation by
\begin{enumerate}
    \item the graph space
    \[
    \GraphSpace{\sigma}
    \equiv
    \ProtoSpace{\sigma}\oplus \DataSpace{\sigma}\oplus \AmbientTarget{\sigma},
    \]
    where \(\ProtoSpace{\sigma}=\CoreState{\sigma}\oplus\GapSpace{\sigma}\);
    \item the coordinate projections
    \[
    \StateProj{\sigma}:\GraphSpace{\sigma}\to \ProtoSpace{\sigma},
    \qquad
    \GraphDataProj{\sigma}:\GraphSpace{\sigma}\to \DataSpace{\sigma},
    \qquad
    \GraphShellProj{\sigma}:\GraphSpace{\sigma}\to \AmbientTarget{\sigma};
    \]
    \item the affine graph plane
    \begin{equation}
    \GraphAffine{\sigma}
    \equiv
    \left\{
    \bigl(
    (m,n),
    \CoreDataMap{\sigma}(m),
    \CoreShellMap{\sigma}(m)
    \bigr):
    (m,n)\in \ProtoSpace{\sigma}
    \right\};
    \label{eq:graphplane}
    \end{equation}
    \item the compact convex graph body
    \begin{equation}
    \GraphBody{\sigma}
    \equiv
    \left\{
    \bigl(
    (m,n),
    \CoreDataMap{\sigma}(m),
    \CoreShellMap{\sigma}(m)
    \bigr):
    (m,n)\in \ProtoBody{\sigma}
    \right\};
    \label{eq:graphbody}
    \end{equation}
    \item the zero-response shell realization
    \begin{equation}
    \GraphEmbed{\sigma}(w)
    \equiv
    \bigl(
    \ProtoEmbed{\sigma}(w),
    \ResponseData{\sigma}(\ProtoEmbed{\sigma}(w)),
    0
    \bigr).
    \label{eq:graphembed}
    \end{equation}
\end{enumerate}
\end{definition}

\begin{proposition}[All graph-level constituents are forced by the completion core]
\label{prop:forcedgraph}
Assume Definitions \ref{def:core}, \ref{def:canonical}, and \ref{def:graph}. Then:
\begin{enumerate}
    \item the affine graph plane in state-plus-response space is necessarily \eqref{eq:graphplane};
    \item the compact convex graph body is necessarily \eqref{eq:graphbody};
    \item the graph-state projection yields the primitive substrate body:
    \[
    \StateProj{\sigma}\bigl(\GraphBody{\sigma}\bigr)=\ProtoBody{\sigma};
    \]
    \item the unique affine graph function is the pair \((\CoreDataMap{\sigma},\CoreShellMap{\sigma})\), and therefore yields the same combined response map by adjoining the hidden-state coordinate:
    \[
    \ResponseMap{\sigma}(m,n)
    =
    \bigl(
    \CoreDataMap{\sigma}(m),
    \CoreShellMap{\sigma}(m),
    n
    \bigr);
    \]
    \item the graph-level hidden splitting is necessarily the state-space factorization
    \[
    \ProtoSpace{\sigma}
    =
    \CoreState{\sigma}\oplus\GapSpace{\sigma},
    \]
    with hidden body \(\GapBody{\sigma}\) and positive operator \(\HiddenOperator{\sigma}\);
    \item the symmetry / time-reversal structure, the exact zero-response shell realization, the fixed-data solvability property, and the graph-kernel coercivity mechanism are all fixed by the same completion core.
\end{enumerate}
\end{proposition}

\begin{proof}
Items (1) and (2) are the definitions of the compressed graph presentation. Because \(\GraphAffine{\sigma}\) is the graph of an affine map over \(\ProtoSpace{\sigma}\), the projection \(\StateProj{\sigma}\big|_{\GraphAffine{\sigma}}\) is an affine isomorphism, which gives item (3). Item (4) is immediate from the graph formula \eqref{eq:graphplane}: once the completion plane is fixed, no second affine graph function is compatible with the same state coordinates.

Item (5) is exactly the product splitting \(\ProtoSpace{\sigma}=\CoreState{\sigma}\oplus\GapSpace{\sigma}\) already used in Definition \ref{def:canonical}. The hidden factor and its gap operator are not extra graph data; they are the same \(\GapBody{\sigma}\) and \(\HiddenOperator{\sigma}\) from the completion core.

For item (6), the induced graph symmetries are obtained by acting on \(\ProtoSpace{\sigma}=\CoreState{\sigma}\oplus\GapSpace{\sigma}\) through the state and hidden components of \(\CoreSym{\sigma}\) and \(\CoreTime{\sigma}\), and on the data and shell coordinates through their visible components. The exact shell realization is \eqref{eq:graphembed}. Fixed-data solvability follows from \eqref{eq:coresolvability}, because the graph data coordinate is determined by the same completion-state coordinate. Finally, if
\[
\delta g=(\delta m,\delta u,\delta e)\in
T_{\GraphEmbed{\sigma}(w)}\GraphAffine{\sigma}
\]
satisfies \(\GraphDataProj{\sigma}(\delta g)=0\), then \(\delta u=0\) and \(\delta m=(\delta m_{\mathrm{co}},\delta n)\) obeys the same data-kernel hypothesis as in Proposition \ref{prop:combined}. Therefore
\[
\|\GraphShellProj{\sigma}(\delta g)\|^{2}
+
\ev{\operatorname{pr}_{\GapSpace{\sigma}}(\delta m),
\HiddenOperator{\sigma}\operatorname{pr}_{\GapSpace{\sigma}}(\delta m)}
\ge
\kappa_{\sigma}^{\mathrm{gr}}\,
\|\delta m\|^{2}
\]
for \(\kappa_{\sigma}^{\mathrm{gr}}=\min\{\kappa_{\sigma}^{\mathrm{co}},\kappa_{\sigma}^{\mathrm{hid}}\}\). Thus the graph-kernel coercivity mechanism is fixed by the same core data.
\end{proof}

\begin{theorem}[Affine response-graph dynamics and combined response law collapse to the completion core]
\label{thm:collapse}
Fix the minimal zero-response affine completion core of Definition \ref{def:core}, together with the canonical combined response law of Definition \ref{def:canonical}. Then:
\begin{enumerate}
    \item the combined response substrate law is canonically induced and necessary by Proposition \ref{prop:combined} and Proposition \ref{prop:forcedcombined};
    \item the affine response-graph presentation is canonically induced and necessary by Definition \ref{def:graph} and Proposition \ref{prop:forcedgraph};
    \item every operative constituent singled out by the previous graph-level note is fixed by the same completion core: the affine graph plane, the compact convex graph body, the graph-state projection yielding the primitive body, the unique affine graph function, the graph-level hidden splitting with positive gap operator, the graph symmetry and time-reversal structure, the exact zero-response shell realization, the fixed-data solvability property, and the graph-kernel coercivity mechanism;
    \item consequently the affine response-graph layer adds no extra benchmark-relevant freedom beyond the completion core and is not a new deepest stopping point on the active line.
\end{enumerate}
\end{theorem}

\begin{proof}
Item (1) is the content of Proposition \ref{prop:combined} and Proposition \ref{prop:forcedcombined}. Item (2) is the content of Definition \ref{def:graph} and Proposition \ref{prop:forcedgraph}. Item (3) is simply the explicit list of forced constituents provided there. Item (4) is the routing consequence: once both the combined-response package and the graph presentation are forced by the same smaller completion core, the graph layer can no longer be treated as an independent remaining freedom.
\end{proof}

\section{Unique Selector and Same One-Metric Output}

\begin{definition}[Deeper data space and sector selector]
\label{def:selector}
For \(\sigma\in\{\Car,\Cur,\Hom\}\), define the deeper data space by
\[
\DeepSpace{\sigma}
\equiv
\operatorname{pr}_{\DataSpace{\sigma}}\bigl(\ResponseImage{\sigma}\bigr).
\]
For \(u\in \DeepSpace{\sigma}\), let \(\ChainSelect{\sigma}(u)\) denote the unique point of \(\ForSpace{\sigma}\) satisfying
\begin{equation}
\ResponseMap{\sigma}\bigl(\ProtoEmbed{\sigma}(\ChainSelect{\sigma}(u))\bigr)
=(u,0,0).
\label{eq:selectordef}
\end{equation}
\end{definition}

\begin{proposition}[The sector selector exists and is unique]
\label{prop:selector}
Assume Definition \ref{def:core}. Then for every \(u\in\DeepSpace{\sigma}\), the point \(\ChainSelect{\sigma}(u)\) of Definition \ref{def:selector} exists and is unique.
\end{proposition}

\begin{proof}
Existence follows from \eqref{eq:combinedsolvability}: if \(u\in\DeepSpace{\sigma}\), then there exists a point of \(\ResponseImage{\sigma}\) with data coordinate \(u\), and \eqref{eq:combinedsolvability} supplies a point of the zero-response shell locus with that same data coordinate. By Proposition \ref{prop:combined}, the map \(\ResponseMap{\sigma}\circ\ProtoEmbed{\sigma}\) is a diffeomorphism onto that zero-response locus, so there is a unique \(w\in\ForSpace{\sigma}\) satisfying \eqref{eq:selectordef}.
\end{proof}

\begin{definition}[Total selector and deeper outputs]
\label{def:deepoutputs}
Define
\[
\DeepShell
\equiv
\DeepSpace{\Car}\times\DeepSpace{\Cur}\times\DeepSpace{\Hom},
\]
\[
\ChainSelectTriple(u_{\Car},u_{\Cur},u_{\Hom})
\equiv
\bigl(
\ChainSelect{\Car}(u_{\Car}),
\ChainSelect{\Cur}(u_{\Cur}),
\ChainSelect{\Hom}(u_{\Hom})
\bigr).
\]
Define
\[
\DeepMetricOut:\DeepShell\to\{\text{observer metrics}\}
\]
and
\[
\DeepMatterOut:\DeepShell\times\Psi\to\{\text{observer stress tensors}\}
\]
by
\begin{equation}
\DeepMetricOut
\equiv
\MetricOut\circ\ChainSelectTriple,
\qquad
\DeepMatterOut(\mathbf u,\psi)
\equiv
\MatterOut(\ChainSelectTriple(\mathbf u),\psi).
\label{eq:deepoutputs}
\end{equation}
\end{definition}

\begin{theorem}[The completion core preserves the same one-metric benchmark package]
\label{thm:outputs}
For every \(\mathbf u\in\DeepShell\) and every matter configuration \(\psi\in\Psi\),
\begin{equation}
\DeepMetricOut(\mathbf u)=\CompletedMetric
\label{eq:deepmetric}
\end{equation}
and
\begin{equation}
\DeepMatterOut(\mathbf u,\psi)
=
T_{\mu\nu}^{\mathrm{matter}}[\CompletedMetric,\psi].
\label{eq:deepmatter}
\end{equation}
Consequently the deeper completion core introduces neither a second operative metric nor an enlarged low-energy matter law.
\end{theorem}

\begin{proof}
Equation \eqref{eq:deepmetric} follows immediately from \eqref{eq:deepoutputs} and \eqref{eq:fixedoutputs}:
\[
\DeepMetricOut(\mathbf u)
=
\MetricOut(\ChainSelectTriple(\mathbf u))
=
\CompletedMetric.
\]
Likewise,
\[
\DeepMatterOut(\mathbf u,\psi)
=
\MatterOut(\ChainSelectTriple(\mathbf u),\psi)
=
T_{\mu\nu}^{\mathrm{matter}}[\CompletedMetric,\psi],
\]
which is \eqref{eq:deepmatter}.
\end{proof}

\section{Main Theorem}

\begin{theorem}[The next remaining burden moves below the graph layer]
\label{thm:main}
Assume the fixed primitive continuation data of Definition \ref{def:fixed} and, for each sector \(\sigma\in\{\Car,\Cur,\Hom\}\), a minimal zero-response affine completion core with hidden-sector gap in the sense of Definition \ref{def:core}. Then:
\begin{enumerate}
    \item the combined response substrate law of Definition \ref{def:combined} is canonically induced by the completion core;
    \item the compact-convex primitive substrate body, the single affine combined response map, the response-image shell / hidden factorization, the hidden-sector gap operator, the induced symmetry / time-reversal structure, the exact zero-response shell realization, the fixed-data solvability property, and the response-kernel coercivity mechanism are all forced by that same core;
    \item the affine graph plane, the compact convex graph body, the graph-state projection yielding the primitive body, the unique affine graph function, the graph-level hidden splitting with positive gap operator, the graph symmetry and time-reversal structure, the exact zero-response shell realization, the fixed-data solvability property, and the graph-kernel coercivity mechanism are likewise all forced by that same core;
    \item the affine response-graph layer and the combined-response layer therefore collapse to the same smaller completion core and differ from any other compatible presentation only by the obvious completion-state, hidden-state, and response-coordinate identifications;
    \item the sectorwise selectors \(\ChainSelect{\sigma}\) and the total selector \(\ChainSelectTriple\) therefore exist uniquely;
    \item pulling the fixed primitive outputs back along that unique total selector yields exactly the same one operative metric \(\CompletedMetric\) and exactly the same ordinary-matter route through that metric;
    \item consequently the remaining burden after the previous affine response-graph theorem is discharged in the stronger form now available on the active line: affine response-graph dynamics are not the new stopping point there, but a redundant presentation of a smaller minimal zero-response affine completion core.
\end{enumerate}
\end{theorem}

\begin{proof}
Item (1) is Proposition \ref{prop:combined}. Item (2) is Proposition \ref{prop:forcedcombined}. Item (3) is Proposition \ref{prop:forcedgraph}. Item (4) is Theorem \ref{thm:collapse}. Item (5) is Proposition \ref{prop:selector} together with Definition \ref{def:deepoutputs}. Item (6) is Theorem \ref{thm:outputs}. Item (7) is the combined routing consequence of the previous items.
\end{proof}

\begin{corollary}[What remains open after this theorem]
\label{cor:next}
After Theorem \ref{thm:main}, the remaining honest burden is no longer to justify the combined response substrate law and it is no longer to justify the affine response-graph layer as an additional deepest package on the active primitive-observer-reduction line. Both have collapsed to a smaller minimal zero-response affine completion core with hidden-sector gap. The next honest burden is therefore deeper again: derive or justify that completion core itself from still more primitive substrate dynamics, or else prove a sharper inevitability theorem that removes even that remaining core-level freedom.
\end{corollary}

\begin{proof}
Theorem \ref{thm:main} converts both the combined-response package and the graph-level package from remaining benchmark-facing burdens into forced presentations of the same smaller completion core. What remains open is why that completion core itself should hold, rather than becoming the current deepest disciplined assumption.
\end{proof}

\section{Conclusion}

The positive result of this note is deliberately narrower than a completed final microphysical derivation and stronger than another deeper same-output relay. The affine response-graph theorem had already shown that the combined response substrate law is forced by a deeper graph package. The present note goes one honest step further by showing that the graph package itself still contains redundant presentation freedom. What actually carries the active line is the smaller minimal zero-response affine completion core: one affine completion plane, one compact convex completion body, one hidden factor with positive gap, the symmetry / time-reversal structure, the exact shell realization, the fixed-data solvability statement, and the core-kernel coercivity inequality.

From that core the note derives canonically the combined response substrate law and the graph-style presentation at once. The compact-convex primitive substrate body, the single affine combined response map, the response-image factorization, the affine graph plane, the compact convex graph body, the graph-state projection, the unique affine graph function, the graph-level hidden splitting with positive gap, the exact zero-response shell realization, the fixed-data solvability property, and the graph-kernel coercivity mechanism are no longer treated as separately inserted primitive ingredients. They are forced presentations of the same smaller completion core.

The benchmark boundary remains fixed throughout. The unique selector determined by the zero-response shell locus still pulls back exactly the same one operative metric \(\CompletedMetric\) and exactly the same ordinary-matter route through that metric, with no second operative metric and no enlarged low-energy matter law. This remains a disciplined derivational gain rather than a final completion claim. The note does not yet derive the completion core itself from ultimate substrate microphysics. Its honest remaining burden is exactly that deeper next step. But the current primary burden named by the project goal is now handled in the stronger form available here: the affine response-graph layer is no longer a new stopping point.

\end{document}
